% from ega2-1-fr.tex
\oldpage[II]{129}

Cela étant, il résulte de (\hyperref[II.6.4.8-fr]{6.4.8}),
(\hyperref[II.6.4.9-fr]{6.4.9}) et (\hyperref[II.6.4.10-fr]{6.4.10}) que si
$s$ est une section de $\mathcal{B}$ au-dessus d'un ouvert $U\subset X$,
$s'$ la section correspondante de $\mathcal{B}'$ au-dessus de $f^{-1}(U)$,
$N_{\mathcal{B}'/\mathcal{A}'}(s')$ est la section de $\mathcal{A}'$ au-dessus
de $f^{-1}(U)$ qui correspond à la section
$N_{\mathcal{B}/\mathcal{A}}(s)$ de $\mathcal{A}$ au-dessus de $U$.

Si $\mathcal{M}$ est un $\mathcal{B}$-Module inversible, on déduit de ce qui
précède que si $\mathcal{M}'=f^*(\mathcal{M})$ (qui est un
$\mathcal{B}'$-Module inversible), on a
$N_{\mathcal{B}'/\mathcal{A}'}(\mathcal{M}')
=f^*(N_{\mathcal{B}/\mathcal{A}}(\mathcal{M}))$ à un isomorphisme canonique
près.
\end{env}

\begin{env}[6.5.5]
\label{II.6.5.5-fr}
Supposons désormais que $(X,\mathcal{A})$ soit un \emph{préschéma}. La donnée
d'une $\mathcal{A}$-Algèbre quasi-cohérente $\mathcal{B}$, qui est un
$\mathcal{A}$-Module \emph{de type fini} équivaut alors, comme on sait, à
celle d'un morphisme \emph{fini} $g:X'\to X$ tel que
$g_*(\mathcal{O}_{X'})=\mathcal{B}$, défini à un $X$-isomorphisme près
(\hyperref[II.6.1.2-fr]{6.1.2} et \hyperref[II.1.3.1-fr]{1.3.1}). En outre,
se donner un $\mathcal{O}_{X'}$-Module quasi-cohérent $\mathcal{F}'$ équivaut
à se donner un $\mathcal{B}$-Module quasi-cohérent $\mathcal{F}$ tel que
$g_*(\mathcal{F}')=\mathcal{F}$ (\hyperref[II.1.4.3-fr]{1.4.3}), et pour que
$\mathcal{F}'$ soit inversible, il faut et il suffit que $\mathcal{F}$ le soit
(\hyperref[II.6.1.12-fr]{6.1.12}). Pour traduire les résultats précédents en
termes du morphisme fini $g$, il faudra donc supposer, soit que
$g_*(\mathcal{O}_{X'})$ est un $\mathcal{O}_X$-Module \emph{localement libre}
(de type fini), soit que $(X,\mathcal{O}_X)$ et $g_*(\mathcal{O}_{X'})$
vérifient l'hypothèse (II). Pour tout $\mathcal{O}_{X'}$-Module inversible
$\mathcal{L}'$, on posera alors
\[
\label{II.6.5.5.1-fr}
N_{X'/X}(\mathcal{L}')
=N_{g_*(\mathcal{O}_{X'})/\mathcal{O}_X}(g_*(\mathcal{L}'))
\tag{6.5.5.1}
\]
et on dira que c'est la \emph{norme} (relative à $g$) de $\mathcal{L}'$. De
même, si $h':\mathcal{L}'_1\to\mathcal{L}'_2$ est un homomorphisme de
$\mathcal{O}_{X'}$-Modules inversibles, on posera
\[
\label{II.6.5.5.2-fr}
N_{X'/X}(h')
=N_{g_*(\mathcal{O}_{X'})/\mathcal{O}_X}(g_*(h')):
N_{X'/X}(\mathcal{L}'_1)\to N_{X'/X}(\mathcal{L}'_2).
\tag{6.5.5.2}
\]
En particulier, pour $\mathcal{L}'_1=\mathcal{O}_{X'}$, on obtient ainsi une
application canonique
\[
\label{II.6.5.5.3-fr}
N_{X'/X}:\Gamma(X',\mathcal{L}')
\to\Gamma(X,N_{X'/X}(\mathcal{L}')).
\tag{6.5.5.3}
\]
Nous laisserons au lecteur la plupart de ces traductions, et nous bornerons à
expliciter les suivantes~:
\end{env}

\begin{proposition}[6.5.6]
\label{II.6.5.6-fr}
Soit $g:X'\to X$ un morphisme fini, et supposons, soit que
$g_*(\mathcal{O}_{X'})$ est un $\mathcal{O}_X$-Module localement libre, soit
que $(X,\mathcal{O}_X)$ et $g_*(\mathcal{O}_{X'})$ vérifient (II \emph{bis})
(ce qui sera le cas en particulier lorsque $X$ est localement noethérien et
normal). Pour qu'un homomorphisme $h':\mathcal{L}'_1\to\mathcal{L}'_2$ de
$\mathcal{O}_{X'}$-Modules inversibles soit un isomorphisme, il faut et il
suffit, dans la première hypothèse, que $N_{X'/X}(h')$ soit un isomorphisme~;
dans la seconde hypothèse, cette condition est encore nécessaire, et elle est
suffisante lorsque l'homomorphisme
$g_*(\mathcal{O}_{X'})\to
g_*(\mathcal{O}_{X'})\otimes_{\mathcal{O}_X}\mathcal{R}(X)$ est injectif.
\end{proposition}

On notera qu'on utilise ici le fait que, pour que $g_*(h')$ soit un
isomorphisme, il faut et il suffit que $h'$ en soit un
(\hyperref[II.1.4.2-fr]{1.4.2}).

\begin{corollary}[6.5.7]
\label{II.6.5.7-fr}
Soit $g:X'\to X$ un morphisme fini, et supposons, soit que
$g_*(\mathcal{O}_{X'})$ est un $\mathcal{O}_X$-Module localement libre, soit
que $(X,\mathcal{O}_X)$ et $g_*(\mathcal{O}_{X'})$ vérifient (II \emph{bis})
et que
$g_*(\mathcal{O}_{X'})\to
g_*(\mathcal{O}_{X'})\otimes_{\mathcal{O}_X}\mathcal{R}(X)$ est injectif.
Soient $\mathcal{L}'$ un $\mathcal{O}_{X'}$-Module inversible, $f'$ une
section de $\mathcal{L}'$ au-dessus de $X'$, $f=N_{X'/X}(f')$ la section de
$\mathcal{L}=N_{X'/X}(\mathcal{L}')$ au-dessus de $X$ qui lui correspond
(\hyperref[II.6.5.5.3-fr]{6.5.5.3}).

